\documentclass[12pt]{article}
\usepackage[utf8]{inputenc}
\usepackage[T1]{fontenc}
\usepackage{hyperref}
\title{Tabaru Language Resources for K-12 Curriculum Development}
\author{Compiled by Claude Sonnet 5 (claude-sonnet-5) for Andr\'as Kornai}
\date{}
\begin{document}
\maketitle

\section{Introduction}
Tabaru is a North Halmahera language of Indonesia with about 15,000 speakers. The search found very little material that could be used in schools. No native-language primer, textbook, or school programme turned up.

\section{Primary Educational Resources}
None were found in search. No Tabaru primers, readers, or textbooks appeared in the results.

\section{Linguistic and Reference Materials}
\begin{itemize}
\item Wikipedia overview of Tabaru. It describes the phonology (five vowels, (C)V syllables, penultimate stress) and cites Kotynski (1988), \emph{Tabaru phonology and morphology}, SIL Work Papers 32, and Fortgens (1928), a Dutch study with folktales: \url{https://en.wikipedia.org/wiki/Tabaru_language}
\item Grambank entry for Tabaru (taba1263), which cites typological work on North Halmahera: \url{https://grambank.clld.org/languages/taba1263}
\item A \emph{Tabaru-Indonesian-English dictionary} (Kotynski, 1994) is mentioned only as a manuscript on a Wiktionary user page. Publication and availability are unverified, and I have no URL I can cite for it.
\item Bowden's \emph{Taba} grammar is about a different language (Taba, South Halmahera). Do not confuse it with Tabaru.
\end{itemize}

\section{Religious and Community Resources}
None were found in search.

\section{Digital and Audio Resources}
None were found in search.

\section{Curriculum Development Context}
Little published material exists for Tabaru. A curriculum would need new documentation and community-led materials. The Kotynski phonology and the 1994 dictionary manuscript would be the starting points, if they can be located.
\end{document}